\begin{otherlanguage}{french}
\chapter*{Résumé}
\addcontentsline{toc}{chapter}{Résumé}
\begin{spacing}{1.5}
Ce projet de fin d'études porte sur la conception et la réalisation d'un écosystème applicatif de communication et de gestion en temps réel nommé \textbf{HMI Stars}. Conçu pour le cabinet de conseil \textbf{HMI Stars Consulting}, cet écosystème vise à simplifier, automatiser et centraliser les échanges administratifs, sociaux et financiers entre le cabinet et ses entreprises clientes. L'architecture repose sur un monorepo Flutter contenant deux applications distinctes connectées à une base de données \textbf{Supabase} : une plateforme web d'administration destinée aux l'entreprise, et une application mobile destinée aux gérants des entreprises clientes.

La plateforme d'administration web permet de superviser les entreprises adhérentes, de collecter les pièces comptables, de planifier les tâches urgentes, de gérer les fiches recus et envoyé et de centraliser les historiques de messagerie. L'application mobile permet aux clients de suivre l'état de leur structure, de déclarer les absences et congés des salariés, de soumettre des pointages d'activités quotidiens, de téléverser des documents administratifs et de dialogue instantané avec le cabinet. Cet écosystème intègre des mécanismes de sécurité via les politiques RLS (Row Level Security) de PostgreSQL et assure la cohérence des flux d'authentification sur le web grâce à la mise en place d'un protocole de routage adaptatif et à l'exploitation du flux d'authentification implicite de Supabase.

\textbf{Mots-clés :} Flutter, Supabase, PostgreSQL, Temps Réel, Messagerie d'Entreprise, Gestion Administrative, RLS.
\end{spacing}
\end{otherlanguage}
